\subsection{Exact order of the vertical-ideal class}

The ideal used in Theorem~\ref{thm:ramified-vertical-truncations}
belongs to a finite chain whose exact sequences determine its order.
This also supplies a split subgroup when the truncation exponent is
not divisible by the square of the residue characteristic.

\begin{theorem}\label{thm:vertical-ideal-exact-order}
Let $O$ be any mixed-characteristic DVR of residue characteristic $p$,
let $\pi$ be a uniformizer, and let $n\geq2$ be divisible by $p$.
Put $s=n/p$, $A=O[t]/t^n$, and $P=C_p$.  Write $X=\sigma-1$ in
$R=A[P]$, and put
\[
 M=(\pi,X+t^s)\subset R,\qquad
 y=[M]-[R]\in K_0(\Lat_A(P)).
\]
Then $y$ and $x=\operatorname{Ind}_P^{S_p}y$ have exact order $p$.
Both lie in the generic and $t$-reduction kernels.  If $p^2\nmid n$,
then each of these classes generates a subgroup $\Z/p$ that is split
in the respective whole lattice Grothendieck group.  For $S_p$ the
retraction can be chosen to annihilate every Young permutation class.
No restriction on the ramification, residue field, or completeness of
$O$ is required.
\end{theorem}

\begin{proof}
Set $Y=X+t^s$ in $R$.  The ring $R$ is free over $A$, with basis
$1,Y,\ldots,Y^{p-1}$, because the change from powers of $X$ is
unitriangular.  Its monic relation in $Y$ has reduction $Y^p$
modulo $\pi$: in residue characteristic $p$ it becomes
\[
 (1+Y-t^s)^p-1=Y^p-t^{ps}=Y^p.
\]
Here $t^{ps}=t^n=0$ in $A$.  Every coefficient below the leading
coefficient of that relation is consequently divisible by $\pi$.

For $0\leq j\leq p$, let
\[
 I_j=(\pi,Y^j)\subset R.
\]
It is free over $A$, with basis
\begin{equation}\label{eq:vertical-ideal-bases}
 \pi Y^q\ (0\leq q<j),\qquad Y^q\ (j\leq q<p).
\end{equation}
To verify this assertion, the indicated span contains $\pi$ and
$Y^j$ and is closed under multiplication by $Y$.  The only last
basis-vector calculation uses the monic relation for $Y^p$; its
coefficients are divisible by $\pi$, as just shown.  Thus the span
is the ideal $I_j$.  Its displayed generators are independent,
since they are diagonal multiples of the basis of $R$ and $\pi$
is a nonzerodivisor on $A$.  The endpoint assertions are
\[
 I_0=R,\qquad I_p=\pi R\cong R.
\]
The second equality again uses $Y^p\in\pi R$.

For $1\leq j<p$ consider the $R$-linear sequence
\begin{equation}\label{eq:vertical-ideal-mesh}
 0\longrightarrow I_j
 \xrightarrow{\ a\mapsto(a,-Ya)\ }
 I_{j-1}\oplus I_{j+1}
 \xrightarrow{\ (b,c)\mapsto Yb+c\ }
 I_j\longrightarrow0.
\end{equation}
Its last map is surjective, since
$YI_{j-1}+I_{j+1}=(\pi,Y^j)=I_j$.
Its first map is injective.  A pair in the kernel of the last map
has the form $(b,-Yb)$, with $b\in I_{j-1}$ and
$Yb\in I_{j+1}$.  Reducing modulo $\pi R$ identifies the ambient
ring with
\[
 (O/\pi)[t,Y]/(t^n,Y^p).
\]
In its free basis $1,Y,\ldots,Y^{p-1}$ over
$(O/\pi)[t]/t^n$, the condition $Y\bar b\in(Y^{j+1})$
forces $\bar b\in(Y^j)$: the only possible extra kernel term
is a multiple of $Y^{p-1}$, which already belongs to $(Y^j)$.
Thus $b\in I_j$, proving exactness in the middle.

The sequence is coefficient-split.  An explicit $A$-linear
section of its last map sends a low basis vector
$\pi Y^q$, $0\leq q<j$, to $(0,\pi Y^q)$, and a high
basis vector $Y^q$, $j\leq q<p$, to $(Y^{q-1},0)$.
These pairs lie in the required ideals and map to their specified
basis vectors.  Therefore \eqref{eq:vertical-ideal-mesh} is an
admissible exact sequence in precisely the lattice category in
question.

Put $z_j=[I_j]-[R]$.  The mesh relations give
\[
 2z_j=z_{j-1}+z_{j+1}\quad(1\leq j<p),\qquad z_0=z_p=0.
\]
Induction on $j$ yields $z_j=jz_1$ and hence $pz_1=0$.
Since $I_1=M$, this says $py=0$.  Exactness of induction gives
$px=0$ as well.

The whole-group detectors in
Theorem~\ref{thm:ramified-vertical-truncations} take $y$ and
$x$ to $s=n/p$ in $\Z/n$.  This value has exact order $p$.
Consequently neither class has smaller order, and both have exact
order $p$.  Their generic and reduction vanishings are the ones
proved in that theorem for this same ideal $M$.

Finally suppose that $p^2\nmid n$, so that $s$ is prime to $p$.
Compose either detector with reduction $\Z/n\to\Z/p$, and
multiply by the inverse of $s$ modulo $p$.  The resulting
homomorphism from the whole Grothendieck group takes the
corresponding class to $1$.  Sending $1\in\Z/p$ to that class
is well-defined because its order is $p$, and is a section of the
homomorphism.  The $S_p$ detector still annihilates the Young span,
so the same is true of this retraction.
\end{proof}

\begin{remark}
When $p^2\mid n$, the homomorphism $\Z/n\to\Z/p$ kills
$n/p$.  The exact-order result still applies, but this composition
does not prove a splitting.  The subgroup generated by this
particular class need not be a direct summand merely because its
order is known.
\end{remark}
